% 本文件由 LaTeX 排版 Agent 根据有效 translation.json 自动生成。
% author=Athanasius; work_id=2807; title=通谕

\chapter{通谕}

% structure: p0001 source=h1 target=omitted reason=page-title-already-rendered-by-parent

\OriginalTerm{通谕信}{Epistola Encyclica}\par

致各地的同为牧职者，至爱的主教们，\Person{亚大纳削}在主内致候安康。\par

% structure: p0004 source=h2 target=section reason=relative-heading-rank; base=h2

\section{§1. 全教会因所发生之事而受影响}

我们所受的苦难，惨烈难忍，非笔墨所能形容；然而，为使你们更易明白所发生之事的可怖性质，我以为最好提醒你们圣经中的一个故事。曾有一名肋未人，其妻子遭受凌辱；当他想到这污辱的极大不洁（因那妇人是希伯来人，属犹大支派），震惊于自己所蒙的暴行，便依照圣经在\BibleBook{民长}{纪}中的记载，将妻子的遗体分割，送到以色列各支派，为使人明了，这等暴行不仅关乎他一人，而是攸关全体；若民众体恤他的苦难，就当为他复仇；若他们置若罔闻，便将从此蒙羞，被视为同有此罪的共犯。他所派遣的使者讲述了所发生的事；凡听闻和目睹的人都说，自以色列子民出\Place{埃及}以来，从未见过如此的事。于是以色列各支派都激愤起来，全体联合攻击那些犯罪者，仿佛自己亲身受苦一样；最终，那恶行的施行者在战争中被消灭，成为众人咒骂的对象：因为聚集的民众不顾念他们乃是同族骨肉，而只看他们所犯的罪。弟兄们，你们熟知这段历史，以及圣经中对此事的详细记载。因此，我不再赘述详情，因为我是写给熟悉这些事的人，并急于向你们的虔敬之心陈述我们当前的境况，这境况比我方才提及的更为恶劣。我提起这段历史的目的，就是要你们将古时的事与现今发生在我们身上的事相比较，看出后者较前者更残酷，从而比古人面对那些罪犯时更加义愤填膺。因为我们所遭受的待遇，比任何迫害都更残酷；肋未人的灾难，与现今对教会所犯下的滔天大罪相比，实在微不足道；更准确地说，此类恶行在全世界闻所未闻，也无人有过类似经历。因为在那个事件中，遭受伤害的只是一个妇女，一个肋未人蒙冤；如今却是整个教会受伤害，司铎职受人凌辱，最可悲的是，虔诚竟遭不虔迫害。那时，各支派因看见一个妇人的尸块而惊骇；如今，整个教会的肢体彼此分离，被送到各处，有的到你们那里，有的到其他人那里，传报他们所受的侮辱和不公。因此，我恳求你们也行动起来，要想到这些冤屈不只临到我们，也同样临到你们；愿各人都伸出援手，感同身受，觉得自己也是受苦者，免得不久之后，教会的法典和信仰都遭败坏。因为若非天主迅速藉你们的手纠正错误，并为教会向她的仇敌复仇，二者都危在旦夕。因为我们的法典与规章，并非今日才赐给各教会，而是由我们的先祖以智慧稳妥地传授给我们。同样，我们的信仰并非此时才起始，而是从主藉着祂的门徒传递给我们的。因此，为使自古保存在教会中直到如今的规章，不致在我们的日子失落，也免致受托于我们的职责在我们手上被追究；弟兄们，你们既是天主奥迹的管家，如今又见这些被他人篡夺，就应奋起。关于我们境况的更多细节，你们将从带信人那里获悉；而我则切愿亲自简要记述，使你们确知：从我们的救主被接升天、命令祂的门徒说：\BibleQuote{你们要去使万民成为门徒，因父及子及圣神之名给他们授洗}以来，教会从未遭受过此等恶行。\par

% structure: p0006 source=h2 target=section reason=relative-heading-rank; base=h2

\section{§2. \Person{额我略}的暴力与非法闯入}

我们所遭受的以及教会所遭受的暴行如下。正当我们如同往常和平集会，信友们在集会中喜乐，在虔敬的交谈中进步，而我们在\Place{埃及}、\Place{底比斯}和\Place{利比亚}的同工们也彼此相爱和睦，并与我们融洽相处时；突然，\Place{埃及}总督发出一封公告，以敕令的形式宣布，一位来自\Place{卡帕多细亚}的名叫\Person{额我略}的，将从宫廷来此接替我。这一宣布使众人错愕，因为这种行径完全前所未闻，此乃破天荒头一遭。然而信友们却更加恒心地聚集在教堂内，因为他们非常清楚，无论他们自己，还是任何主教或司铎，简言之，没有任何人曾控告过我；他们看到只有亚略派人士支持他，并且也知道他本人就是亚略派，是由\Person{优西比乌}及其同党派遣到亚略派来的。弟兄们，你们知道，\Person{优西比乌}及其同党一向是亚略派狂徒们不虔异端的支持者和同伙，他们藉此不断图谋害我，也是我被流放到\Place{高卢}的始作俑者。\par

因此，民众理所当然地愤慨并大声抗议这一行为，呼唤其余的官员和全城的人作证，这桩新奇的、不义的企图如今是针对教会发动的，并非因为教会人士对我提出任何指控，而是由于亚略异端者的恶意攻击。因为即使真的有任何普遍流传的针对我的投诉，也不该选择一个亚略派的人，或一个宣认亚略教义的人来取代我；而应按照教会的法典，以及\Person{保禄}的指示，当民众\BibleQuote{聚集在一起，而授圣职者的心意}，并具\BibleQuote{我们的主\Person{耶稣基督}的权能}时，一切都应在要求更换的平信徒及圣职人员面前，依照法典进行查问和办理；而不该由亚略派从远方带来一个人，仿佛在买卖主教名衔，且凭借异教官吏的庇护和强力，强行闯入那些既不邀请也不渴望他出现、甚至对所发生之事一无所知的人中间。这类行径导致所有教会法典的破坏，迫使异教徒亵渎，并怀疑我们的任命不是依照神圣的规则，而是出于买卖和庇护的结果。\par

% structure: p0009 source=h2 target=section reason=relative-heading-rank; base=h2

\section{§3. \Person{额我略}到来时所发生的暴行}

\Person{额我略}的这次显赫任命就是这样由亚略派促成的，而这就是其开端。至于他进入\Place{亚历山大}时所犯下的暴行，以及这件事成了多少巨大祸患的原因，你们既可从我们的书信中得知，也可通过询问寄居在你们中间的人了解。当民众对这种不寻常的做法感到愤慨，因而聚集在教堂里，为阻止亚略派的不虔敬与教会的信仰相混合时，\Person{斐拉格留斯}——他长期以来就是教会及其贞女的迫害者，如今是\Place{埃及}总督，业已背教，且是\Person{额我略}的同乡，一个品性卑劣之人，更得到\Person{犹西比乌斯}及其同伙的支持，因此对教会满怀狂热；此人通过后来兑现的承诺，成功拉拢了异教群众，连同犹太人和不法之徒，煽动他们的情绪，用刀剑棍棒将他们成群地派往教堂攻击民众。\par

随后发生的事绝非轻易能够描述：事实上，不可能向你们公正地呈现这些情形，甚至若不流泪哀叹，连一小部分都无法诉说。这类行为在古代可曾被作为悲剧题材？或者在迫害或战争的时期，曾有类似之事发生过吗？教堂和\OriginalTerm{圣洗池}{baptistery}被点燃，随即全城充斥着呻吟、尖叫和哀号；市民对这些暴行愤慨不已，斥责总督，抗议施加给他们的暴力。因为圣洁无玷的贞女被剥去衣服，遭受难以启齿的对待，倘若抵抗，便有性命之虞。隐修士被践踏致死；有的被猛力抛下；有的被刀剑和棍棒杀害；有的被打伤、殴打。啊！在\OriginalTerm{祭台}{Holy Table}上犯下了何等不虔不义的行径！他们用鸟和松果献祭，歌颂他们的偶像，甚至就在教堂里亵渎我们的主、救主、生活天主之子\Person{耶稣基督}。他们在教堂中找到的圣经书卷便付之一炬；而杀害我们主的凶手——犹太人，以及邪恶的异教徒，竟无礼地（哦，何其大胆！）闯入圣洗池，脱去衣服，言行举止如此可耻，人甚至羞于述说。还有某些不虔之人，效法最残酷迫害中的榜样，抓住贞女和苦修士的手，拖着他们走，边拖边强迫他们亵渎并否认主；当他们拒绝时，就残暴地殴打并践踏他们。\par

% structure: p0012 source=h2 target=section reason=relative-heading-rank; base=h2

\section{§4. 339年圣周五与复活主日的暴行}

此外，在这样显赫而辉煌地进入城市之后，\OriginalTerm{亚略派}{Arian}的\Person{额我略}对这些灾祸幸灾乐祸，并似乎渴望给予\OriginalTerm{外教人}{heathens}、\Person{犹太人}，以及那些对我们施加这些恶行的人，作为他们不义得逞的奖赏和报酬，就把这座圣堂交出，任凭他们抢掠。因着这罪恶与混乱的特许，他们的所作所为比战时更糟，比强盗更残忍。有的抢夺一切到手的东西；有的瓜分某些人存放那里的钱财；大量的酒，或喝，或倒掉，或抬走；掠夺油的储藏，每人都把门和圣所栏杆当作掠物拿走；蜡台他们立刻靠墙放好，并在他们的偶像前点燃圣堂的蜡烛：总之，抢掠和死亡弥漫了圣堂。而那些不虔的\OriginalTerm{亚略派}{Arian}，非但不为这些事感到羞耻，反而变本加厉，行出更多暴行与残忍。\OriginalTerm{司铎}{presbyters}和平信徒被撕碎皮肉，\OriginalTerm{贞女}{virgins}被扯去头纱、押解到总督的法庭，然后投入监狱；另一些人财产被充公，并遭鞭打；\OriginalTerm{圣职人员}{ministers}和贞女的食物被截断。这些事竟发生在\OriginalTerm{四旬期}{Lent}圣季，在\OriginalTerm{复活节}{Easter}前后；那时弟兄们正守斋戒，而这位显赫的\Person{额我略}却展现了\Person{盖法}的心性，与总督\Person{比拉多}一同，狂暴地攻击基督的虔诚信众。在\OriginalTerm{预备日}{Preparation}那天，他同总督和一群外教人进入一座圣堂，见到众人对他的强行闯入憎恶不已，便指使那个最残忍的人，即总督，在一个小时内公开鞭打三十四名贞女、已婚妇女和显贵人士，并将他们投进监狱。其中有一位贞女，热爱学习，在她被下令公开鞭打时，正手持\WorkTitle{圣咏集}；书被差役撕成碎片，那位贞女本人被关进监牢。\par

% structure: p0014 source=h2 target=section reason=relative-heading-rank; base=h2

\section{§5. \Person{亚大纳削}的隐退，以及\Person{额我略}与\Person{菲拉格里乌斯}的暴政}

当这一切发生之后，他们仍不罢休，反而商议如何在另一座圣堂行同样的事，那些日子我大多住在那里；他们迫不及待地要将怒火延烧到那座圣堂，以便搜出我并置我于死地。若非基督的\OriginalTerm{恩宠}{grace}护佑我，这原本就是我的命运——我之所以逃脱，也许只为能将这些关於他们行径的少许细节讲述出来。因为我见他们对我狂怒至极，且忧虑圣堂受损，里面的\OriginalTerm{贞女}{virgins}受害，更多凶杀发生，百姓再次受凌辱，便从他们中间退去，记起我们救主的话：\BibleQuote{若有人在这一城迫害你们，你们就逃往另一城去！}我知道，从他们对第一座圣堂所作恶事来看，他们对另一座圣堂也必不会手下留情。事实上，在那里，他们连圣节期的主日也不尊重，反而在那座圣堂里也拘禁了属於它的人——那时，正是主解救众人脱离死亡的锁链之际，而\Person{额我略}及其同党，仿佛在与我们的救主争战，又倚仗总督的庇护，竟将基督之仆获得自由的日子变成了哀悼。外教人欢喜行此事，因为他们憎恶那个日子；\Person{额我略}或许只是执行\Person{欧瑟伯}及其同党的命令，强迫基督徒在捆锁的痛苦中哀悼。\par

透过这些暴力行径，总督夺取了各座圣堂，将它们交给了\Person{额我略}和\OriginalTerm{亚略派}{Arian}的狂徒。如此一来，那些因不虔而被我们开除教籍的人，如今竟以掠夺我们的圣堂为荣；而天主的子民和\OriginalTerm{天主教会}{Catholic Church}的\OriginalTerm{圣职人员}{Clergy}，要么被迫与\OriginalTerm{亚略异端者}{Arian heretics}的不虔\OriginalTerm{共融}{communion}，要么就被禁止进入圣堂。此外，藉着总督，\Person{额我略}还向船主和其他过海的人施以不少暴力，拷打鞭笞一些人，给另一些人套上锁链并投入监狱，以迫使他们不反抗他的恶行，并接受他发出的信件。这一切他尚不满足，为了大口喝我们的血，他还让那个野蛮的同党——总督，以百姓的名义，向最虔敬的皇帝\Person{君士坦提乌斯}控告了我；其捏造的罪状十分可憎，人们据此可以预料到的，不只是被放逐，甚至是死千万次。起草这罪状的人，是一个背弃基督信仰、肆无忌惮敬拜偶像的背教者；签名支持它的，有外教人、看守偶像庙宇的人，还有一些是亚略派。简言之，为避免我的信件对你过于冗长，我要说：这里正掀起一场\OriginalTerm{教难}{persecution}，一场前所未闻的对教会的迫害。因为从前，人们至少可以在躲避迫害者时祈祷，在藏匿时受\OriginalTerm{圣洗}{baptism}。但如今，他们极度的残忍竟效仿了巴比伦人渎神的行为。因为正如他们诬告\Person{达尼尔}，这位显赫的\Person{额我略}如今也在总督面前控告那些在家中祈祷的人，并伺机凌辱他们的圣职人员；因此，由于他的暴行，许多人因错失圣洗而陷入危险，许多患病忧伤的人无人探望；这种不幸使他们痛哭哀号，认为比患病本身更糟糕。因为当教会的圣职人员遭受迫害时，那些谴责亚略异端者不虔的子民，宁可带病冒险，也不愿让亚略派的手落在他们头上。\par

% structure: p0017 source=h2 target=section reason=relative-heading-rank; base=h2

\section{§6. 上述一切非法行为皆是为\OriginalTerm{亚略派}{Arianism}的利益而进行}

\Person{额我略}是\OriginalTerm{亚略派}{Arian}，已被派往亚略派团体；因为无人要他，只有他们；所以像佣工和外人一样，他利用总督向天主教会的人民施行这些可怕而残忍的事，好像这些人不是他自己的。因为先前\Person{欧瑟伯}及其同党所立管理亚略派的\Person{比斯都}，因着他的不虔敬，已由你们天主教会的主教们公正地绝罚和开除教籍，正如你们所知，在我们写信告知你们关于他的事后，他们现在以同样方式把\Person{额我略}派到他们那里；并且为了避免因我们再次写信反对他们而再次受辱，他们用外来的力量攻击我，以便占据教会，似乎可以逃脱亚略派的嫌疑。但在这点上他们也错了，因为教会的人民中除了异端者，以及因各种罪被开除教籍的人，和被总督强迫假装的人之外，无人与他们同在。这就是\Person{欧瑟伯}及其同党长期排练和导演的戏剧；现在他们通过向皇帝对我做虚假控告，成功地演出了。然而他们还不满足于安静，甚至现在还想杀害我；他们使我们的朋友们极其恐惧，以致他们都被放逐，并预期死在他们手中。但你们不要因此害怕他们的邪恶，反而要报复：对这种对我们前所未有的行为表示愤怒。因为\BibleQuote{若是一个肢体受苦，所有的肢体都一同受苦}\BibleRef{格前 12:26}，并且按照真福宗徒的话，\BibleQuote{与哭泣的一同哭泣}\BibleRef{罗 12:15}，那么如今这样大的一个教会受苦，每一个人都应如同自己受苦一样为她伸冤。因为我们有一位共同的救主，他们亵渎了他，还有属于我们大家的\OriginalTerm{法典}{Canons}，他们正在违反。如果你们中任何一位正坐在你们的教会里，人民与你们聚集，没有任何过错，突然有人假借敕令，作为你们中一位的继任者而来，并对你们做了同样的事，你们岂不愤怒？岂不要求伸冤？如果这样，那么现在你们理应愤怒，以免这些事被忽视，同样的祸患逐渐蔓延到每个教会，我们的宗教学校变成市场和交易所。\par

% structure: p0019 source=h2 target=section reason=relative-heading-rank; base=h2

\section{呼吁全教会的主教联合反对\Person{额我略}}

你们熟知\OriginalTerm{亚略狂人}{Arian madmen}的历史，亲爱的，因为你们常常个人地或集体地谴责他们的不虔敬；你们也知道\Person{欧瑟伯}及其同党，如我以前说的，从事同样的异端；为了这异端，他们长期对我进行阴谋。我已向你们说明了现在他们已经并正在以比战时更残酷的方式所作的事，为的是让你们以我在开始时所述历史中的例子为前导，对他们的邪恶满怀热忱的憎恨，并拒绝那些对教会犯下如此暴行的人。如果在这些事发生之前，\Place{罗马}的弟兄们[去年]因着他们以前的恶行，写信要召集一次\OriginalTerm{大公会议}{Council}，以便纠正这些恶事（\Person{欧瑟伯}及其同党担心此事，事先设法扰乱教会，想要毁灭我，以便此后他们可以随心所欲地行事而不怕，没有人能追究他们），那么，你们现在该是多么更应该对这些暴行表示愤怒，并谴责他们，因为他们在前恶之上又加了这些。\par

我恳求你们，不要忽视这些事，也不要让著名的\Place{亚历山大里亚}教会受到异端者的践踏。由于这些事，人民和他们的圣职人员彼此分离，正如所料，因\OriginalTerm{长官}{Prefect}的暴力而沉默，却憎恶亚略狂人的不虔敬。因此，如果\Person{额我略}或任何其他代表他的人写信给你们，弟兄们，不要接受他的信，要撕碎它们，并使带信的人蒙羞，因为他们是邪恶和不虔的仆役。即使他假装友善地写信给你们，也不要接受。那些带信的人只是由于害怕总督以及他频繁的暴行才送信。并且，\Person{欧瑟伯}及其同党很可能就他的事写信给你们，因此我急切地预先警告你们，好使你们在此效法那不看情面的天主，把从他们那里来的人从你们面前赶出去；因为他们为了亚略狂人的缘故，致使外教人和犹太人在这样的时节，在教会里进行迫害、强奸贞女、谋杀、掠夺教会财产、纵火和亵渎。那邪恶而疯狂的\Person{额我略}无法否认他是亚略派，写他信的人证明了这一点。这人就是他的秘书\Person{阿孟}，他很久以前因许多恶行和不虔敬被我前任真福亚历山大逐出教会。\par

因着这一切原因，恳请你们给我回信，并谴责这些不虔敬的人；好使此地的圣职人员和人民，看到你们的正统信仰和对邪恶的憎恨，因你们在基督徒信仰上的同心而喜乐，也使那些对教会犯了这些不法行为的人因你们的信而改过，最终虽迟却悔改。请问候你们中间的众弟兄。所有与我在一起的弟兄都问候你们。你们要安好，记念我，愿主常保护你们，至亲爱的众位主。\par

\SourceRecord{Translated by M. Atkinson and Archibald Robertson.；Nicene and Post-Nicene Fathers, Second Series；Vol. 4.；Edited by Philip Schaff and Henry Wace.；Buffalo, NY: Christian Literature Publishing Co.,；1892.；<http://www.newadvent.org/fathers/2807.htm>.}
